\section{Discussion}
The clearest result is that uncertainty recognition can vary independently of evidence-following. Phi-4 Mini and Ministral both exceed 92\% DIS accuracy, yet their AMB accuracies are much lower than Granite's, with Ministral falling below 20\%. A model can therefore answer well when the context is explicit while still making unsupported person-level judgments when the correct behavior is to abstain.

The results also show why directional bias scores cannot be interpreted alone. Ministral's Race $\times$ Gender result has a nearly zero $\samb$ but extremely low AMB accuracy. Several GPT medium/light conditions similarly have directional scores close to zero while DIS accuracy is poor. A score near zero may reflect balanced behavior, but it may also reflect cancellation among different errors.

Category-level reporting remains important because aggregate values hide concentrated effects. Physical Appearance and Age repeatedly produce larger ambiguous-context bias for the local models, while Religion and Physical Appearance appear among stronger disambiguated hotspots. These findings should still be interpreted within BBQ's U.S. English-speaking cultural scope \citep{parrish2022bbq}.

The observed negative alignment gaps suggest that some conditions perform worse when the correct answer conflicts with the stereotype-aligned target. However, the overall gaps are small and the larger category-level gaps are based on subsets, so we treat them as descriptive rather than statistically significant.

Finally, the GPT conditions require additional caution. Their answer sequences can be rescored, but the original inference procedure is unavailable. The two high conditions are especially audit-limited because their sequences are identical and perfectly match the gold answers. We therefore use them only transparently and avoid strong comparative claims.
